.mt9
.mb7
.po8
.ll65
.ig
CHAPTER 1
U S I N G   T H E   D O C U M E N T A T I O N
.sp2
.pp5
Before you operate the Payroll System, read this manual thoroughly
and do the Practice Run outlined in Chapter 4.
.he Introduction: Chapter 1
.pp
The SSG Payroll System documentation includes two manuals.  This
manual is a Reference Manual which provides all the background information
and instructions necessary for successful operation.  The other is a
Specifications Manual which contains technical information on file accessing
and file descriptions.	The Specifications Manual is available only
with the purchase of the system itself.
.pp
The Reference Manual contains an Introduction, a Reference Section, an
Operations Section, and an Appendix.  The Reference section deals with
general concepts and procedures.  The Operations section gives
additional specifics on running the programs. After you get a good
understanding of the system from the Reference section, you
should rarely have to refer to it.  From then on, most questions can
be answered by referring to the Operations section.
.pp
While Applewood Computers is glad to provide support to
registered owners of its products, please
read the documentation carefully and check for hardware
problems before giving us a call.
.he
.bp
.ig
CHAPTER 2
F E A T U R E S    A N D    F A C I L I T I E S
.sp2
.pp5
The Applewood Computers Payroll System is a sophisticated payroll
processing program for small and medium businesses.  A two drive
single-density 8" system can handle about 50 employees paid monthly
and submitting daily time cards.  The same microcomputer system can
handle approximately 100 employees if only one or two pay transactions
are needed each pay period.  A double-density system can handle about
200-300 such employees.  Systems with more than two drives will
be able to handle more employees.
.he Introduction: Chapter 2
.pp
The Payroll System calculates deductions from employee wages
for Federal income tax withholding, Social Security tax (FICA), State Disablity
Insurance (SDI), and state and local income tax withholding for many state
and local governments.	The Payroll System also calculates the employer's
FICA, FUTA (Federal Unemployment Tax), and SUI (State Unemployment
Insurance) contributions.  Advance Earned Income Credit (EIC) payments
and reduction in payments for earnings over the EIC limit are also
handled by the system.
.pp
In addition to the standard deductions listed above, you can define
up to five system-wide earnings or deductions and up to three employee
specific earnings or deductions.  Of course any employee can be exempted
from the standard and system-wide earnings and deductions categories.  User
defined earnings and deductions can be a flat amount each pay period,
or a percentage of the employee's gross taxable pay.  You can limit
the amount that can be earned or deducted in one year, or the amount
earned or deducted can be without limit.  The user-defined earnings
categories can be marked as subject to all taxes and percentage based
deductions, subject to all except FICA, subject to all except Federal,
state, and local taxes, or not subject to any taxes or percentage based
deductions.
.pp
Provision has been made for entering an employee's head of household
status and special withholding allowances for California tax purposes.
.pp
The Payroll System supports additional Federal, state, local, and FICA
withholding if requested by an employee.
.pp
The system handles tips collected by the employer, and tips
reported to but not collected by the employer.	The Payroll System
exempts the employer from FICA tax on tips, and taxes the employee
correctly for them.
.pp
The Payroll System permits three fully taxable pay rates, plus one
rate with a tax status that may be defined differently for different employees.
These normally
correspond to regular pay, overtime pay, special overtime, and
commission, piecework,pay, or other pay, though the descriptive name may be any
you desire.
.pp
The system can generate recurring pay amounts automatically (with a
simple override procedure), and calculate overtime pay automatically.
The Payroll System calculates vacation and sick leave
earned, and provides for comp time entry.
.pp
The Payroll System handles employee bonuses,
expense reimbursement, salary advances and repayments, sick pay, and
vacation pay.
.pp
It supports weekly, bi-weekly, monthly, and semi-monthly pay schedules for
employees classified as active, on leave, or terminated.
.pp
It will print paychecks automatically, and if the option is requested,
void checks over a stated maximum amount.  As an added safety feature,
the Payroll System gives an error message when an employee
receives payment in
excess of the maximum you set for him or her.
.pp
You can distribute an employee's total cost to the company to up
to five ledger accounts.  The system is completely
auditable, and interfaces
with the Structued Systems Group General Ledger.
.pp
The Payroll System prints 941 government reporting
forms every quarter, and 940 and W-2 forms
at the end of every year.
.pp
It produces payroll transaction working and audit proof reports, a
Payroll Journal that details earnings and deduction information by
employee (with a company wide and ledger account summary),
and a paycheck register from
which checks may be prepared manually if desired.  The system prints
out the master employee file in a full or short form for some
or all employees, and also prints an employee history report
which presents summarized quarter and year-to-date pay information.
You can order a report of vacation, sick leave, and comp time
accumulated and used, as well as
a printout of the accounts currently used by the Payroll System.
.pp
Structured Systems' Display Management System provides full screen
formatting to make data entry a simple matter of filling in the blanks.
Default values for all requests speed data entry and
ensure successful processing.  The Applewood Computers Payroll
System incorporates extensive processing safeguards and error
checking routines, and employs our most advanced defensive programming
techniques.
.he
.bp
.ig
CHAPTER 3
S Y S T E M    O V E R V I E W
.sp2
.pp5
The Applewood Computers Payroll System is "menu driven", which
means you run the programs by choosing them from a "menu" list.
The Payroll System employs SSG's Display Management System, a full-screen
display formatting utility that enables you to
view a number of requests and their current values at once.  To enter
a new value or change an old one, you move the cursor to the desired
field and type in the new response.  The Payroll System edits your
responses to its requests and checks them carefully for reasonableness
and validity, issuing descriptive error messages when necessary.
.he Introduction: Chapter 3
.pp
The Payroll System breaks logically into three parts: System Setup,
the Payroll Processing Cycle, and End of Period Processing.  Once
installation of your Payroll System is complete, you will rarely need
to take the System Setup steps again.  The Payroll
Processing Cycle is repeated once for every pay period, and End of Period
Processing once at the end of each quarter and year.
.sp2
.ce
SYSTEM SETUP
.pp
Step One of System Setup is to make two copies of the "distribution"
disks you received from SSG.  One set you use only for making
additional copies, the other set is for actual processing.
.pp
Step Two of System Setup is to describe the characteristics of you
CRT (console) device to the Display Management System.	For most
this is a simple renaming process. Others may need to run the
CRT definition program described in Chapter 33.
.pp
Step Three of System Setup involves erasing the unneeded state tax
table files from disk, and renaming the tax table file
for your local government.
.pp
Step Four of System Setup is called Parameter Entry.  During this
step you tell the system what options you want to use, as well as
other information to make processing more efficient.  At first time
startup the Payroll System lets you choose between using a simpler
"default" version of the system (Quick Parameter Entry),
and doing further customization.
If you choose the default Payroll System, the parameter requests you
answer during system startup are limited in number.
If you choose further customization at first time startup,
the system presents you with the full range of parameter requests, and
allows you to change account and deduction information.
If all employees are paid on the same schedule, and if no additional
accounts or special earnings or deduction categories are needed, the
default system will be the simplest setup procedure.  Of course, you can
always choose to do further customization later.
.pp
Step Five of System Setup is Account Entry.  If you choose the
default Payroll System option at first time startup, no requests are
issued and the program creates a default account file containing
a cash, accrued liablity, salary expense, and accrued payroll
cost liability account.
If you choose further customization, you can add additional accounts,
or change the existing default accounts.  If you select the General
Ledger interface option, the Account Entry program prompts you to
insert the disk with your GL chart of accounts file so it can
verify the account information you enter.
.pp
Step Six of System Setup is called System-wide Earnings and Deduction
Entry.	During this step you tell the system which of the standard
deduction categories you wish to use (such as Federal, state, and
local income tax, FICA, SDI, FUTA, EIC, etc.), and enter current
rate and limit information.  During this step you can define up to
five system-wide earnings or deductions to suit the specific needs of your
company.  Individuals may of course be exempted.  If you choose
the default Payroll System option, at first time
startup the program creates a default Earnings and Deductions file and
allows no entries or changes.  You can, however, make changes to the
default file at a later time through the Deduction/Earning Entry program.
.pp
Step Seven of System Setup is called Employee Entry.  There are two
programs to be run for this step.  The first is Employee Entry,
which creates a record on the Employee Master and History
files for each
employee, and assigns him or her a number.  The second is the
Employee Deduction/Earnings and Cost Distribution
program, which you can use to
exempt specific employees from system-wide earnings or deductions,
establish up to three earings or deductions categories for a particular
employee, and distribute an employee's cost to the company to up
to five ledger accounts.
.pp
Step Eight of System Setup is Employee Pay Rate Entry.	During this step
you enter an employee's pay frequency, set the four
pay rates (usually regular, overtime, special overtime, and
commission, although you can give them any description), and
set the number of Normal Units (used for automatic computation of
overtime), and set the number of Autogen Units (used for automatic generation
of recurring pay amounts).  During this step you also enter the
rate at which vacation and sick leave are accumulated, and the
vacation, sick leave, and comp time accumulated and used so far in the
current year.
.pp
When you set up your Payroll System in mid-year, a ninth setup step is
required.  After the first eight steps are complete, you
run the History Entry program to bring the earnings, deduction,
and tax information up to date for each employee and for the company
as a whole.
It will not run after the first
payroll application has taken place.
See Chapter 7 and Chapter 19 for instructions.
.sp2
.ce
PROCESSING CYCLE
.pp
When setup is complete, you can begin the processing cycle
you repeat for each pay interval to calculate employee earnings and deductions,
produce paychecks, and update historical information.  The Payroll System
permits four pay intervals: weekly, bi-weekly, monthly, and semi-monthly.
The weekly and bi-weekly schedules are called week-based intervals, and
the monthly and semi-monthly schedules are called month-based intervals.
Although pay transactions for all four intervals can be entered into
the same pay transaction batch, the central program of the processing
cycle, Apply Payroll, never processes week-based and month-based
employees together.  When weekly and bi-weekly employees are
on the system, the Apply Payroll program processes them together
every two weeks.  When monthly and semi-monthly employees are on the system,
the Apply Payroll program processes them together at the end of
each month.
.pp
For the purpose of this overview, the Payroll Processing Cycle can be
broken into four steps: enter pay transactions, sort the transaction file,
apply payroll for the period, and print paychecks and/or the paycheck
register.
.pp
The first step in the cycle, Enter Pay Transactions, is for entering
employee time, vacation or sick leave used, comp time earned, tips
reported or collected, bonuses, advances on pay, advance repayment,
and other transactions.  A working proof report of the transactions entered
so far is available at any time.  When transactions have been entered for
all the employees to be paid by the next payroll application, a final proof
for the audit trail must be generated.
.pp
After making backup copies of your disks, you take the second step in
the cycle, Sort the Transaction file.
.pp
The third step in the cycle is to Apply Payroll for the period.  The
Apply Payroll program is the main processing program of the system.
It separates
transactions for week-based employees from those for month-based,
and if both intervals for the interval base are not to be processed
by this application (e.g., semi-monthly, but not monthly), it separates
those that are to be processed from those that are not.  The Apply Payroll
program produces a check file from which the paychecks and
the Paycheck Register are produced, updates the Company History and
Employee History files, and creates a General Ledger batch file if
that option has been requested.  Before you can run the application program
again, you must
print a Payroll Journal report that presents detailed and summarized
transaction information for each employee, a company payroll summary,
and a ledger account summary.
.pp
The fourth step in the cycle is to actually print the pay checks and
the Paycheck Register.	If you do not choose the automatic check
printing option, you can use the Paycheck Register to manually prepare
the paychecks.
.sp2
.ce
END OF PERIOD PROCESSING
.pp
After the last payroll application of each quarter,
and before the first payroll application
of the next quarter, you must print the Federal 941 report
(Employer's Quarterly
Federal Tax Return), and run the
Close for Quarter program to reset quarter to date statistical
registers to zero.
.pp
After Closing for the fourth quarter, and before the first payroll application
of the next year, you must print the Federal 940 report,
print W-2 forms for all employees,
and run the Close for the Year program to reset
year to date statistical registers to zero.
.he
.bp
.ig
CHAPTER 4
P R A C T I C E    R U N
.sp2
.pp5
Before you begin this practice run, follow the procedure outlined in
Chapter 19 for making copies of the SSG distribution disks, installing
the Display Management System, and setting up the tax tables for your state
and local government. It is assumed you are using state tax tables.  If you
are not you must run the Deduction/Earning Entry program and set
State Withholding Used to No before running
the Apply Payroll program.
Follow the Practice Run
instructions carefully, making the
entries exactly as directed.  When you are done,
you may want to work with the practice run data for
several processing cycles more before starting for real.
.he Introduction: Chapter 4
.pp
Make sure you have the programs, tax files, and if needed, the CRT.DEF file
on the proper disks.  If you are using the three disk configuration (Startup,
Processing, and Data disk), you will need to replace the startup disk with
the processing disk after you have completed Employee Rate Entry.
If you are using the two disk configuration (a Program disk, and a Data disk),
you will not need to change disks.
See that CRUN2.COM (the CBASIC run time interpreter
program) is on the currently logged disk drive, and make sure
the disk you put in drive A has an operating system on it. Hit the RESET switch
on your machine, and
when the CP/M prompt character is showing, call up the Payroll System by
typing:
.sp
.li
		A>CRUN2 PY
.sp
After you press RETURN, the CRUN version number, the SSG copyright message,
the Payroll System version number, and
the serial number of your Payroll Package appear.
.pp
When the Payroll System menu appears select number 10, Quick Parameter Entry.
Be sure to hit RETURN after each response.  If you key in the
wrong character, give the command CTRL-H to back up and erase one character
(hold down the CONTROL key and type the letter "H").
.pp
When the Parameter Entry screen appears, hit RETURN until the cursor is
at the Company Name field.  If you go by, you can back up
by giving the command CTRL-B.
Type in the information
below, hitting RETURN after each entry:
.sp
.li
   NAME:       Universal Practice Company
   ADDRESS 1:  2120 Main Street
   ADDRESS 2:  [hit RETURN]
   ADDRESS 3:  [hit RETURN]
   CITY:       Anytown
   STATE:      [enter the 2 character abbreviation of your state]
   ZIP:        60609
   PHONE:      (408)222-3322
   GOVT ID NUMBERS:
     FEDERAL:  01-0092688
     STATE:    72-7228809
     LOCAL:   [hit RETURN to skip the local ID number]
.sp
Make no other entries or changes.  Hit ESCAPE to end Parameter Entry.
Answer the date request with 1/16/80.
After RETURN,
the system creates an account file, a deduction file, and
a sort parameter file, and the menu returns.  Choose number 24, Employee
Entry.
.pp
When the Employee Entry screen appears, type "RUN" to create the Employee
Master and History files.
When the cursor is waiting at the
Employee Number field, give the command:
.sp
.li
		CTRL-A
.sp
The program will display the first available employee number and
move the cursor to the NAME field.  Type in the following information
for the first employee:
.sp
.li
   NAME:      Brian Boisterous
   ADDRESS 1: 747 Airport Way
   ADDRESS 2: Apartment 3B
   CITY:      Anytown
   STATE:     [your state]
   ZIP:       60715
   PHONE:     (408)635-2929
   SOCIAL SECURITY NUMBER:     167-82-4906
   BANK ACCOUNT NUMBER:        8705-05654
   BIRTH DATE:		       11/21/46
   SEX:        M
   PENSION:    N
   JOB CODE:   [hit RETURN]
   START DATE: 7/4/75
   EMPLOYMENT STATUS: [hit RETURN]
   TERM. DATE: [hit RETURN]
.sp
Now press ESCAPE to begin the next employee.
Remember to press RETURN after each response.  You can go back one
field by CTRL-B, forward one field by RETURN, and erase an erroneous
character by CTRL-H.  Enter the information below for the second employee.
.sp
.li
   NAME:      Margaret Mellow
   ADDRESS 1: 23 Lovely Lane
   ADDRESS 2: [hit RETURN]
   CITY:      Anytown
   STATE:     [your state]
   ZIP:       60715
   PHONE:     (408)635-3030
   SOCIAL SECURITY NUMBER:     916-22-1066
   BANK ACCOUNT NUMBER:        893-25619
   BIRTH DATE:		       4/1/51
   SEX:        F
   PENSION:    Y
   JOB CODE:   [hit RETURN]
   START DATE: 12/15/74
   EMPLOYMENT STATUS:  [hit RETURN]
   TERM. DATE: [hit RETURN]
.sp
Press ESCAPE when the entries are correct.  Hit ESCAPE again
to leave the add-records mode, and then a third time
to end the Employee Entry program.  When the menu returns select number
26, Employee Pay Rate Entry (we will skip the Employee Ded/Earn and Cost Dist.
program since the default values are fine for now).  When the Employee Pay
Rate Entry screen appears with the cursor waiting at the Employee Number
field, give the command:
.sp
.li
		CTRL-N
.sp
to bring up the employee record for the first employee you entered.
Press RETURN to begin entering information for the employee.
Leave the Pay Type set to Salaried (S) and the Pay Interval set to
Semi-Monthly (S) by hitting RETURN at those requests.  Enter the following
rate information (do not enter commas, do enter the decimal):
.sp
.li
   REGULAR:	   1150.75
   OVERTIME:	   [hit RETURN]
   SPEC. OVERTIME: [hit RETURN]
   COMMISSION:	   1
.sp
Answer the remaining requests with the following information:
.sp
.li
   RATE4 TAX EXCLUS. TYPE: [hit RETURN]
   AUTOGEN UNITS:	   1
   NORMAL UNITS:	   [hit RETURN]
   MAXIMUM PAY: 	   5000
   MARITAL STATUS:	   [hit RETURN]
   FEDERAL WH ALLOWANCES:  1
   STATE   WH ALLOWANCES:  1
   LOCAL   WH ALLOWANCES:  [hit RETURN]
   EIC USED (Y OR N):	   [hit RETURN]
.sp
Leave the remaining values unchanged and hit ESCAPE to begin the
next employee.	When the cursor is at the Employee Number field,
give the CTRL-N command to get the next employee record.  After the
name and social security number appear, press RETURN to begin entering
the information below
.sp
.li
   PAY TYPE:		   H
   PAY INTERVAL:	   [hit RETURN]
   REGULAR:		   5
   OVERTIME:		   7.50
   SPECIAL OVERTIME:	   10.00
   COMMISSION:		   [hit RETURN]
   RATE4 TAX EXCLUS. TYPE: [hit RETURN]
   AUTOGEN UNITS:	   [hit RETURN]
   NORMAL UNITS:	   80
   MAXIMUM PAY: 	   [hit RETURN]
   MARITAL STATUS:	   M
   FEDERAL WH ALLOWANCES:  3
   STATE   WH ALLOWANCES:  3
   LOCAL   WH ALLOWANCES:  [hit RETURN]
   EIC USED (Y OR N):	   [hit RETURN]
.sp
Make no more entries and press ESCAPE to end the employee.  Press
ESCAPE once more to end the Employee Rate Entry program.  The setup steps
are now complete. If you are using the three disk configuration, press ESCAPE
once more to end the Payroll System.  Replace the Startup Disk in A with the
Processing disk, give the command CTRL-C, and then call up the Payroll System
by typing CRUN2 PY 1/16/80.
.pp
Assume that the first working period (January 1 to January 15) is now
over.  Today is January 16. Today you will enter pay transactions, generate
an audit proof report of the transactions, sort the batch of pay
transactions, apply payroll for the period, print paychecks (on plain paper),
and print the Paycheck Register.  To begin, select menu item number 1,
Pay Transaction Entry.
When the Pay Transaction Entry screen appears, hit RETURN at the Display
Employee Names request.  When
the cursor is waiting at the Record Number field, give the command:
.sp
.li
		CTRL-A
.sp
to begin adding transactions to the file.  Type in the transaction information
below, hitting ESCAPE to begin the next transaction. Do not enter the Record
Number, it is presented for reference only.
.sp
.cp12
.po11
.li
Record	Employee   Transaction		  Effective
Number	 Number       Code	 Units	    Date

  1	 00014	       8	   8	   [RETURN]  [ESCAPE]
  2	[RETURN]      10	  43.55    [RETURN]  [ESCAPE]
  3	[RETURN]      11	 200	   [RETURN]  [ESCAPE]
  4	[RETURN]       4	 124.75    [RETURN]  [ESCAPE]
  5	 00028	       0	  88	   [RETURN]  [ESCAPE]
  6	[RETURN]       9	  50.00    [RETURN]  [ESCAPE]
  7	[RETURN]      15	  13.12    [RETURN]  [ESCAPE]
  8	[RETURN]      18	  22.12    [RETURN]  [ESCAPE]
  9	[RETURN]      28	  25.12    [RETURN]  [ESCAPE]
.po8
.sp
After entering in the last transaction, press ESCAPE to leave the add-records
mode,
and ESCAPE again to end the Pay Transaction
Entry program and return to the menu.
.pp
Align the paper in your printer (132 column printer is assumed), choose
menu item number 2, Pay Transaction Proof.
If the paper should jam, hitting any key
on the keyboard will end the printout.
.pp
When the printout is finished and the menu returns, choose menu item
number 3, Pay Transaction Sort.  When the sort is complete the
menu will return.
.pp
Choose menu item number 4, Apply Payroll, and when asked, type "YES"to
confirm that the interval to be processed is correct.
Give the Check Date as 1/16/80.  When asked, type "YES" to create the
Company History file.
The menu returns when processing is over.
.pp
Make sure the paper in your printer is aligned correctly, then print
the Payroll Journal by choosing menu item number 5, Print Payroll
Journal.
.pp
Now select the Print Paychecks program, menu item number 6.  Use
plain paper. Answer N (NO) to the test form request.
Give the starting check number as 1 and then press RETURN.
.pp
When the menu returns, choose item number 7, Print Check Register.
.pp
When the menu returns after the Check Register is printed,
select number 20, Print Employee Master List.  When the printing option
menu appears, select number 1, ALL EMPLOYEES.  Enter F at the request
for the Full or Single line option.  When the printout ends and the print
option menu returns, press ESCAPE to return to the main system menu.
.pp
Press ESCAPE at the system menu to end the Payroll System.
Continue practicing with this data, or erase
all files except the INT files, the COM files,
the tax files, and the CRT.DEF file
(if used) before starting for real.  When the CP/M prompt appears, remove
the disks and store them safely away.
.he
.bp
.ig
CHAPTER 5
G A T H E R I N G    Y O U R	M A T E R I A L S
.sp2
.pp5
Having the information you need right at hand makes system setup an
easy task.  After you read the documentation thoroughly, decide what
Payroll System options you want to use and gather the
data you'll need for each employee and for the company as a whole.
.he Introduction: Chapter 5
.pp
Many of the requests in the Parameter Entry program are for "default"
values.  Defaults are items of information used in the absence of other
instructions (i.e., "in default").  The default values you set in
Parameter Entry appear as the current values in the corresponding
Employee Entry requests. In certain situations they can save a great
deal of time.  For example, if many employees accumulate vacation at
the same rate, you would enter that rate at the request for the default
vacation rate, and thus avoid having to key in the same answer for
many employees.  If you think carefully about the default values you
enter, you might be able to save yourself a substantial amount
of time and effort.  Specifically,
the Parameter Entry program lets you set default values for:
.sp
.cp7
.li
	1. Vacation Rate
	2. Sick Leave Rate
	3. Pay Rates 1 through 4
	4. Rate4 Tax Exclusion Type
	5. Employee Pay Interval (e.g., weekly, monthly, etc.)
	6. Pay Type (hourly or salaried)
	7. Maximum Pay Factor (see Chapter 8)
	8. Normal Units (see Chapter 8)
.pp
The Payroll System account file created automatically by the system contains
four default accounts: cash, accrued liabilities, salary expense, and accrued
payroll costs liability.  If you plan to add or change accounts, or if you've
chosen the GL interface option, it will save time to have your chart
of accounts handy also.
.pp
By far the best thing you can do to speed system setup along is to have the
personnel information you need gathered in one place before you sit
down to the machine.  In addition to having the information you need for
setting up the system-wide earnings and deductions (such as FICA rate
and limit, EIC rates, etc.), you'll want to have information on earnings
and deductions that apply to specific employees (such as insurance or
savings account deductions).
.pp
An Employee Data Collection Form is provided (Figure 5.1) to help you
gather personnel data in a form that requires little or no adaptation
to be entered into the Payroll System.	We suggest you make as many
copies of this form as you need (plus a few extra), and fill one out
for each employee.  If you're unsure how to answer any request, re-read
the documentation for a full explanation.
.he
.bp
.mt4
.mb5
.ce2
Applewood Computers PAYROLL SYSTEM
EMPLOYEE ENTRY DATA COLLECTION FORM
.sp2
.li
Name:_____________________________________________________

Address:__________________________________________________

	__________________________________________________

City:____________________  State:_____	Zipcode:__________

Phone: (      ) ________-______________

Social Security No.: ______-____-________

Bank Account No. (optional) ______________________________

Birthdate: ____/____/____     Sex: ______

Pension Plan:	 Y   or  N   (circle one)

Job Code (3 characters, optional):  __________

Starting Date: ____/____/____

Employment Status: Active  On Leave  Terminated  (circle one)

Termination Date (if applicable): ____/____/____

Distribution of Employee Cost to Company:

	      Account No.	     Percentage

	      __________	     __________

	      __________	     __________

	      __________	     __________

	      __________	     __________

	      __________	     __________
			    Total:	 100 %

Exempt From Withholding?

	      Federal	     Y	  or   N     (circle one)
	      FICA	     Y	  or   N     (circle one)
	      State	     Y	  or   N     (circle one)
	      SDI	     Y	  or   N     (circle one)
	      Local	     Y	  or   N     (circle one)
	      FUTA	     Y	  or   N     (circle one)
	      SUI	     Y	  or   N     (circle one)
.sp
.ce
(Figure 5.1)
.bp
.li
Pay Type:    Hourly   Salaried	   (circle one)

Pay Interval:  Semi-Monthly  Monthly  Weekly  Bi-Weekly  (circle one)

Rate1 Pay Rate:  _______________

Rate2 Pay Rate:  _______________

Rate3 Pay Rate:  _______________

Rate4 Pay Rate:  _______________  Rate4 Tax Exclusion Type:  ______

Automatic Pay Units (see Chapter 11):  ____________

Normal Pay Units (see Chapter 11):  ____________

Maximum Pay (see Chapter 11):  ____________

Marital Status:   Married    Single    (circle one)

Withholding Allowances:   Federal __________

			  State   __________

			  Local   __________

Advance Earned Income Credit Option Used?   Y	or   N	(circle one)

State of California Special Information:
	   Head of Household:	 Y   or   N   (circle one)
	   Special Withholding Allowances: __________

Vacation Rate:	_______________

Vacation Accumulated This Year:  _______________

Vacation Used This Year:  _______________

Sick Leave Rate:  _______________

Sick Leave Accumulated This Year:  _______________

Sick Leave Used This Year:  _______________

Comp Time Accumulated This Year:  _______________

Comp Time Used This Year:  _______________
.sp
.ce
(Figure 5.1)
.he Introduction: Chapter 5
.mt7
.mb9
.bp
.ce
MID-YEAR SETUP
.pp5
To bring the Payroll System up to date when starting in mid-year, you
run the History Update program (see Chapters 7 and 19).  This program
requests quarter and year to date earnings and deduction information
for each employee, and for the company as a whole.  You also enter
quarter and year to date information on company FICA, FUTA, and SUI
costs, and quarter to date vacation earned and taken.  Finally, the
program asks you to enter the amount and date of Federal
tax deposits for the current quarter, and the amount of Federal,
FICA, and FUTA tax payments for previous quarters.  See Chapter 32
for reproductions of the entry screens.  An Employee History Update Data
Collection Form is provided for your convenience.  Make as many copies
of this form as needed and fill one out for each employee (Active,
On Leave, and Terminated).
.he
.mt5
.mb5
.bp
.ce3
Applewood Computers PAYROLL SYSTEM
EMPLOYEE HISTORY UPDATE DATA COLLECTION FORM
Q U A R T E R	 T O	D A T E    I N F O R M A T I O N
.sp2
.li
Employee Name: ______________________________  Number: ____________



	  QUARTER TO DATE EARNINGS BY TAX CATEGORIES

1. Income (subject to all taxes)		   ________________

2. Other (subject to all but FWT, SWT, and LWT)    ________________

3. No Taxes (income not subject to taxation)	   ________________

4. FICA (income subject to FICA tax; not	   ________________
   including tips)
5. EIC CREDIT (advance EIC payment received)	   ________________

6. TIPS (collected and reported)		   ________________

7. NET (sum of checks written)			   ________________



		QUARTER TO DATE TAXES WITHHELD

       FEDERAL __________   STATE __________  LOCAL __________

		 FICA __________  SDI __________



	     QUARTER TO DATE EARNINGS AND DEDUCTIONS

SYS1 ________ SYS2 ________ SYS3 ________ SYS4 ________ SYS5 ________

	  EMP1 ___________  EMP2 ____________  EMP3 ____________

OTHER DEDUCTIONS: _______________

UNITS BY PAY TYPE: 1 _________	2 _________  3 _________  4 _________
.sp3
.ce
(Figure 5.2)
.bp
.ce3
Applewood Computers PAYROLL SYSTEM
EMPLOYEE HISTORY UPDATE DATA COLLECTION FORM
Y E A R    T O	  D A T E    I N F O R M A T I O N
.sp2
.li
Employee Name: ______________________________  Number: ____________



	    YEAR TO DATE EARNINGS BY TAX CATEGORIES

1. Income (subject to all taxes)		   ________________

2. Other (subject to all but FWT, SWT, and LWT)    ________________

3. No Taxes (income not subject to taxation)	   ________________

4. FICA (income subject to FICA tax; not	   ________________
   including tips)
5. EIC CREDIT (advance EIC payment received)	   ________________

6. TIPS (collected and reported)		   ________________

7. NET (sum of checks written)			   ________________



		  YEAR TO DATE TAXES WITHHELD

       FEDERAL __________   STATE __________  LOCAL __________

		 FICA __________  SDI __________



	      YEAR TO DATE EARNINGS AND DEDUCTIONS

SYS1 ________ SYS2 ________ SYS3 ________ SYS4 ________ SYS5 ________

	  EMP1 ___________  EMP2 ____________  EMP3 ____________

OTHER DEDUCTIONS: _______________

UNITS BY PAY TYPE: 1 _________	2 _________  3 _________  4 _________
.sp3
.ce
(Figure 5.2)
.he
.mt7
.mb9
.bp
.ig
CHAPTER 6
B A C K G R O U N D   I N F O R M A T I O N
HARDWARE AND SOFTWARE REQUIREMENTS
.sp2
.pp5
The SSG Payroll System
will run on any 8080 or Z-80 based microcomputer
under the control of the CP/M operating system or a CP/M compatible
operating system.  Your microcomputer system must have the following
minimum hardware and software components:
.he Introduction: Chapter 6
.pp
1.  54K (kilobytes) of Random Access Memory (RAM)  (CP/M configured)
.pp
2.  Two floppy disk drives with a minimum capacity of 256K each,
although additional disk capacity is recommended.
.pp
3.  A 132 column wide printer capable of skipping to the top of the next
form in response to the ASCII form-feed character.
.pp
4.  For the console device, a CRT (Cathode Ray Tube) with a direct cursor
addressing facility, a clear screen control, and
minimum dimensions
of 24 lines by 80 columns.
.pp
5.  System software including an operating system and the CBASIC2 language
(version 2.04 or earlier).
.pp
Many microcomputers can be expanded or adapted to meet these requirements.
Your local dealer can give you specific information about the hardware
and software currently available.
.sp2
.cp6
.ce
THE CBASIC PROGRAMMING LANGUAGE
.pp
The programs that make up the Payroll System are written
in a dialect of the BASIC programming language, called CBASIC2.
CBASIC2 is a product of Software Systems; questions relating to
the CBASIC2 language should be referred to Software Systems.
.pp
Applewood Computers distributes its programs in an intermediate
file format (with the filetype INT), which must be further translated
at run time by a
facility of the CBASIC2 language called the "interpreter".  The interpreter's
program name CRUN2 must precede any INT program for it to run.
For example, to run the Payroll System you would enter CRUN2 PY.
This run time interpreter comes as a part of the CBASIC2
language package.
With the CBASIC2 language you
can also write programs yourself to extend the capabilities
of the Payroll
System, or to perform other data processing tasks.
.sp2
.cp6
.ce
THE CP/M OPERATING SYSTEM
.pp
An operating system is a set of instructions which enables the computer to,
among other things, read and write information on floppy disks, move and
maintain files, and provide for console communication.	The Payroll
System
employs the CP/M operating system, or one compatible with CP/M.  CP/M (TM) is a
product of the Digital Research Company. Questions concerning the
CP/M system specifically should be referred directly to that company.
.pp
CP/M comes documented by a series of six manuals.  The first of the six,
"Features and Facilities", contains useful and important information
relating to your operation of the Payroll System.  Be particularly
careful to read and understand the DIR, ERA, REN, PIP, SYSGEN, TYPE,
and STAT commands before using the Payroll System.
.sp2
.cp6
.ce
LOGGING ON
.pp
When you first turn your computer on, its random access memory (RAM) is
empty.	Since in this state the computer can do nothing,
a feature has been built into the microcomputer hardware that
directs it to look on a predetermined section of a certain disk drive to find
what to do next.  If everything is working normally, your computer should
look on the outer two tracks of the A disk, find the CP/M operating system
(if it has been SYSGEN'ed onto the disk; see the CP/M "Features and Facilities"
manual for details on how this is done), load it into the RAM, and then
display "A>" on the CRT, indicating it is ready to go to work.  This
process is called "booting" since the computer is, in effect, lifting itself
up by its own bootstraps.
.pp
The "A>" character is called a "prompt".  When the A> prompt appears on the
screen, the operating system is said to be "logged on" to the A disk drive.
While logged on to the A drive, you can call up programs or data files on that
disk by typing the program name.
To request a file on another disk (unless this is done automatically
by the program you are using) first log on to that disk by keying in
the drive name (the letter B on a two disk drive system), followed by a
colon (:) and a carriage return, and then type the program name.
(The carriage return is simply the RETURN
key as it is on a typewriter; it is sometimes referred to as CR.)
.sp
.cp10
.li
	For example,
.sp
.li
	   A>		   [system prompt]
	   A>B: 	   [user enters B:, then RETURN]
	   B>
.sp
To return to the A drive, enter:
.sp
.li
	   B>A: 	   [user enters A:, then RETURN]
	   A>
.pp
Whenever the prompt appears you are operating under the CP/M
environment.  It goes away when you enter the Payroll System, and
returns when you leave it.
.sp2
.ce
FULL-DISK BACKUP
.pp
Caution and common sense dictate that whenever you purchase or receive a
new program or system  on a floppy disk, you immediately copy it
onto a fresh disk and store the original in a fire-proof, moisture-proof  box
to protect it against inevitable human error and unpredictable system failure.
With the original SSG distribution disk tucked safely away from the ravages of
dirt, grease, magnet, and mangle, you avoid the cost and inconvenience of
unnecessary replacement.  This warning applies doubly for disks
containing data crucial
to your operation.  Do not make the mistake of being "penny wise and pound
foolish"; the disk itself is almost always worth far less than the information
residing on it.
.pp
The manufacturer of your disk drives should have provided a full-disk copy
program with the disk hardware.  If this program is not available, use the
PIP command that comes with the CP/M operating system. To do so,
place a properly formatted disk containing an
operating system and the PIP program on
drive A, and the Payroll System
distribution disk (the disk you received from Applewood Computers)
on drive B.  Hit the RESET
switch on your machine and then type
.sp
.li
	A>PIP A:=B:*.*[OV]
.sp
This will copy all the programs and files on B onto the A disk and verify
that they were copied correctly.
See the CP/M documentation
for a complete description of the PIP program.
.sp2
.ce
CHANGING DISKS BETWEEN PROGRAMS
.pp
Whenever you change disks between programs it is of the utmost importance
that you enter a CONTROL-C character immediately after inserting the new disk.
The
CTRL-C character, which appears on your screen as ^C, informs the computer
that a new disk with different data has been inserted.	This
"re-booting" causes the CP/M system to reset its internal
maps of disk usage, and prevents it from writing over (thus destroying)
files which exist on the new disk.
.pp
To enter CTRL-C, or any CONTROL character, hold down the CONTROL key (CTRL
or CTL on some machines), while at the same time pressing the letter C
(either upper or lower case).
.he
.bp
.ig
CHAPTER 7
S Y S T E M    S E T U P
.sp2
.pp5
This chapter covers the initial steps in setting up your Payroll System:
copying the distribution disks, setting up the
Display Management System, and setting up the tax tables for your
state and locality.  It also explains the sequence of events that occurs
when the Payroll System is run for the first time, and the mid-year startup
procedure.
Chapter 8 deals with setting the Payroll System
parameters.  Chapter 9 covers setting up the Payroll System accounts
Chapter 10 covers establishing the earnings or deductions that
are to apply to all employees (unless exempted), and
changing the Federal Withholding Tax Table.  Chapter 11 treats the
final setup steps, putting employees on the system, and entering their
pay rate information.
.he Reference: Chapter 7
.pp
The Payroll System is distributed on three diskettes.  One disk holds
the programs needed for setting up your Payroll System (Startup Disk).
One holds the
programs needed for normal processing (Processing Disk).
And one holds the state and local
tax table files, the CRT definition files, the CRT definition program
CRTDEF.INT, and the single density copy program SDCOPY.COM (see the
Payroll System Specifications Manual).	This third distribution disk
is called the Tax Disk, but the copy you make of that disk will be called
the Data Disk once you begin.
.pp
The disks you receive are called the distribution disks.  They should
never be used for normal processing.  Two sets of copies should be made
of the distribution disks, and they should then be stored safely away.
One set should be labelled the "System Master Disks" and used to make
additional copies of the programs that may be needed.  The other set
should be labelled "Work Disks" and may be used for actual processing.
(See Figure 7.1).
.pp
If you are using double density disks (or otherwise have at least
500K per disk), you need only two disks to run the Payroll System, a program
disk and a data disk.
One disk that size will hold all the programs (INT and COM files), plus
PIP, STAT, and CRUN2.  If you are operating in single density, you will
need three disks, a startup disk, a normal processing disk, and a data disk.
Once CRUN2 is on the disks with the programs, PIP and STAT will most likely
not fit on the disk also.  There should be no trouble fitting them on the
data disk, however.
.sp
.cp34
.li



						     DISTRIBUTION DISKS


      STARTUP DISK     PROCESSING DISK	   TAX DISK



						     SYSTEM MASTER DISKS



      STARTUP MASTER   PROCESSING MASTER  TAX MASTER



							  WORK DISKS
						     (3 DISK CONFIGURATION)


      STARTUP WORK   PROCESSING WORK	PAYROLL DATA
	 DISK		 DISK		    DISK


							  WORK DISKS
						     (2 DISK CONFIGURATION)



      PAYROLL PROGRAM DISK	  PAYROLL DATA DISK

		 (Figure 7.1: Disk Labelling)
.pp
If you are using the three disk configuration, you will need to change disks
when the setup steps are complete (from the Startup disk to the Processing
disk), and whenever you want to run a program that is not on the disk currently
inserted.  The menu program is distributed on both the Startup and the
Processing disks.  CRUN2 and CRT.DEF (if used) should be on both disks.
.cp5
.sp
STARTUP DISK MUST BE IN DRIVE A FOR:
.sp
.in6
.ti-3
1. The Payroll System first time startup sequence.
.sp
.ti-3
2. Adding, changing, or examining employee records (i.e., using Employee
Entry, Employee Ded/Earn and Cost Dist., and Employee Pay Rate Entry).
.sp
.ti-3
3. Adding, changing, or examining accounts (i.e., using Account Entry).
.sp
.ti-3
4. Printing the Account List.
.sp
.ti-3
5. Changing or examining Standard deduction information, the Federal
Tax Table, and System-wide earnings and deductions.
.sp
.ti-3
6. Changing or examining system parameters.
.sp
.ti-3
7. Running the History Entry program.
.qi
.sp2
.cp5
.li
THE PROCESSING DISK MUST BE IN DRIVE A FOR:
.sp
.in6
.ti-3
1. Pay Transaction Entry, Proof, and Sort.
.sp
.ti-3
2. All print programs except Print the Account List.
.sp
.ti-3
3. Payroll Application.
.sp
.ti-3
4. End of Period Processing.
.qi
.pp
After copying the distribution disks, you install the Display
Management System.  Because CRT devices differ widely, the Display
Management System (DMS) must know the characteristics of your CRT.
The Display Management System comes ready to run on a Hazeltine 1500 CRT
device.
CRT definition (DEF) files are provided on the distribution disks for
a number of other popular CRT's listed in Chapter 19.  If you
have a CRT for which there is a CRT definition file, a simple
renaming process is all that's needed. If not, you must run the program
described in Chapter 33 to create one before you continue.
.pp
The next step is to set up the state and, if needed, local tax tables.
Applewood Computers distributes the Payroll System with a number
of state and local tax table files.  Each file is named to indicate
which state or local government it is for.  Use the CP/M DIR command
to locate the file for your state and erase the other state files.  To
set up the local tax table, you locate the one for your local government
and rename it according to the instructions in Chapter 19.
.pp
Normally when you call up the Payroll System a menu appears from which
you choose the programs to run.  However, the first time the
system is run you are automatically taken
through a sequence of startup steps before the full range of menu options
is made available to you.
After you give the command to call
up the Payroll System (CRUN2 PY), the CRUN version number appears
at the top of your screen, followed by the SSG copyright message,
the Payroll System version number, and the serial number of your Payroll System
package.
The system menu then appears with a message prompting you to choose
one of the two parameter entry programs (Quick Parameter Entry or Parameter
Entry),
depending on whether you want to use
the default parameter values, accounts, and systemwide deductions,
or customize the system further by answering
additional parameter requests and entering other account and deduction
information.
Read the chapter on Setting System
Parameters (Chapter 8) and note the default values presented beside
each request to decide which startup option to choose.	Read also the
chapters on the Account File (Chapter 9) and System-wide Earnings and
Deductions (Chapter 10), noting the defaults created by the programs
described in those chapters.
.pp
If you choose
the default Payroll System (Quick Parameter Entry),
after you answer the limited number
of parameter requests, the system goes on to create the default
account file, the default system-wide deduction and earnings file,
and the sort parameter file.  Later, if you need to, you can change
the default values in these files by choosing the appropriate program
from the menu and making the entries in the normal way.  (Default values
are values used in the absence of other instructions, i.e., "in default".)
.pp
If you do not
use the default Payroll System, after you answer all the parameter requests,
the system goes on to the Account Entry program, and then the System
Deduction Entry program where you can make entries or changes.
After these steps, it creates the sort parameter file used by the sorting
program (there are no requests), and then presents you with the system
menu, from which you can run the final setup programs.
.sp2
.cp6
.ce
MID-YEAR SETUP PROCEDURE
.pp5
The accuracy of many of the Payroll System reports (especially the
government reporting forms), and the correctness of certain earning
and deduction calculations (for example FICA, which applies only to
a limited part of an employee's earnings), depend on historical
information you enter through the History Entry program distributed
with the Payroll System.
.pp
The History Entry program should be run after the Employee Pay
Rate Entry setup step is complete.  It can be run as many times as
necessary before the first run of the Apply Payroll program.  After
that point, the program is automatically disabled to protect the
audit trail.  The program is run from the menu.
.pp
Quarter to date and year to date earnings, withholding, and deduction
information must be entered for each employee, active, on leave, and
terminated.  The Employee History Update Data Collection Form
in Chapter 5 is provided for your convenience.	Reproductions of the
entry screens can be found in Chapter 32.  The program requests the
following information for each employee, totals it, and then presents
the company-wide totals for modification.  The year to date requests
have been omitted since they
are identical to the quarter to date requests.
.sp
QUARTER TO DATE EARNINGS BY TAX CATEGORIES:
.sp
.in8
.ti-8
INCOME: This request is for income
subject to all taxes earned
in the current quarter (FWT, SWT, LWT, FICA, FUTA, SUI, SDI).
.sp
.ti-8
OTHER:	This request is for income subject to all except Federal, state, and
local income taxes
earned in the current quarter (FICA, SDI, FUTA, SUI).
.sp
.ti-8
NO TAXES: This request is for income not subject to taxation received
in the current quarter.
.sp
.ti-8
FICA:	This request is for income subject to FICA tax earned in the
current quarter. Do not include tips since they are handled separately.
.sp
.ti-8
EIC CREDIT: This request is for the amount of advance Earned Income
Credit payments received so far in the current quarter.
.sp
.ti-8
TIPS:	This request is for the amount of tips collected (by the employer)
and reported (by the employee to the employer) so far in the current quarter.
.sp
.ti-8
NET:	This request is for the net pay received so far in the current
quarter.  This is normally the sum of checks written to the employee.
.qi
.sp
QUARTER TO DATE TAXES WITHHELD
.sp
.in8
.ti-8
FEDERAL: Enter the amount of Federal Taxes withheld from the employee
so far in the current quarter.
.sp
.ti-8
STATE:	Enter the amount of state taxes withheld from the employee so far
in the current quarter.
.sp
.ti-8
LOCAL:	Enter the amount of local taxes withheld from the employee so far
in the current quarter.
.qi
.sp
QUARTER TO DATE EARNINGS AND DEDUCTIONS
.sp
.in8
.ti-8
SYS1 THROUGH SYS5: If you have established systemwide earnings or deduction
categories through the Deduction/Earning Entry program, these requests
let you enter the amounts paid or withheld so far in the current quarter
for each employee.
.sp
.ti-8
EMP1 THROUGH EMP3: If you have established employee specific earnings or
deduction categories through the Ded/Earn and Cost Distribution
program, these
requests allow you to enter the amounts paid or withheld so far in the
current quarter for each employee.
.sp
.ti-8
OTHER DEDUCTIONS: If you have made any other deductions from an employee's
wages enter the amount deducted so far in the current quarter.
.sp
.ti-8
UNITS BY PAY TYPE 1 THROUGH 4: Enter the number of pay units each
employee has accumulated so far in the current quarter.  For example,
if an employee has worked 120 regular hours and 20 overtime hours, you
would enter 120 for Rate (Type) 1, and 20 for Rate2. The information you
enter has no effect on calculations or government reports.  It is
for informational purposes only.
.qi
.sp2
.pp
After you answer the requests described above for quarter and year
to date information for each employee, the program updates each record
on the Employee History file. It then totals the information
to arrive at a company-wide total, which it presents for modification when
you choose Update Employee History Totals from the "mini" menu that appears
when you are done entering information for each employee.
.pp
A second menu choice, Update Company Liabilities, Vacation, Etc., lets
you enter the amount of the FICA, FUTA, and SUI liabilities incurred
by the company (NOT withheld from wages)
in the current quarter and year.  Since the vacation,
sick leave, and comp time information you entered through the Employee
Rate Entry program is for the year to date, this entry screen lets you enter
company-wide quarter to date information if desired.  The year to date
vacation, sick leave, and comp time information is totaled from each
employee record and presented also.  You can modify this information
if necessary.
.pp
A third menu choice, Update Payments History, lets you enter the amount
and date of the Federal tax payments you have made so far in the
current quarter.  The total Federal, FICA, and FUTA tax payments you have
made in each of the quarters so far during the year are also requested.
.he
.bp
.ig
CHAPTER 8
SETTING SYSTEM PARAMETERS
.sp2
.pp5
Parameters are items of information that tell the Payroll System
what options you want to exercise, where it can find certain files it
needs during processing, and certain other things it needs to know to operate
efficiently.  Because not all users have the same needs, you can adapt
the Payroll System to your requirements by changing parameter (pronounced
pe-'ram-et-er) values,
rather than by expensive and complex reprogramming.
.he Reference: Chapter 8
.pp
Payroll System users fall into one of two classes.  Some can
use the default
parameter values that come with the Payroll System, and will only need
to answer a few parameter requests (for their company name, ID
numbers, etc.). Others may want to
customize the system further by answering additional parameter requests
and by changing account and deduction information.
To decide which class you belong to, read the description of each
parameter request
presented in this chapter and note whether the default value
is adequate for your needs.  If you choose the default Payroll system
at startup (Quick Parameter Entry), you
need only answer the first set of requests presented below
(Parameter Requests that Must Be Answered In Any Case).
Whether you use the "default" Payroll System,
or choose to customize it further, the first set of parameter requests
shown below must be answered in any case.  Of course, if your needs change
you can change the parameter values later through the regular Parameter
Entry program.
.sp
.ce
PARAMETER REQUESTS THAT MUST BE ANSWERED IN ANY CASE
.in5
.sp
.ti-5
PAY INTERVAL USED [S]: You must use at least one of the four possible
pay schedules: weekly (W), bi-weekly (B), monthly (M), or semi-monthly (S).
The default Payroll System allows only one pay interval.  If you need more,
you should choose the option to do further customization.
.sp
.ti-5
DATE FORM: [MM/DD/YY]: State whether you want to enter dates in
the American (MM/DD/YY)
or international date format (DD/MM/YY).
.sp
.ti-5
ENDING DAY OF LAST PERIOD [12/31/79]: If only one pay interval is used,
you would enter the last day of the last working period.  The Payroll
System adds one day to arrive at the first day of the new period, and
the appropriate number of days to arrive at the last day. The ending day
for the monthly interval must be the last day of the month.  The ending
day for the semi-monthly interval must be the 15th of the month, or the
last day of the month.	The last day of the last period for weekly or
bi-weekly intervals can be any you desire.
.sp
.ti-5
LAST QUARTER ENDED [4]: State which quarter was the last quarter ended.
.sp
.ti-5
CURRENT YEAR [1980]: Enter the current year.
.qi
.sp
.li
FEDERAL ID NUMBER:
STATE	ID NUMBER:
LOCAL	ID NUMBER:
.in5
.sp
.ti-5
LINES PER PAGE [ 66]: Enter the length in lines. The print programs
assume 132 column wide paper.
.sp
.ti-5
DEBUGGING AID [ No]: Leave set to No.
.sp
.ti-5
CONSOLE BELL [Yes]: When set to Yes, the console alarm sounds if
an error occurs.
.sp
.ti-5
COMPANY NAME: This field allows 60 characters, but since the W-2 forms
permit only 30 characters and the
941 form accepts only 43 characters, an abbreviated version may be
preferable.  If the company name exceeds these limits, the excess is
truncated.
.qi
.sp
.li
ADDRESS: Three lines are provided for the street address.
CITY:
STATE:
ZIP:
PHONE:
.sp2
.ce2
ADDITIONAL REQUESTS YOU ANSWER
IF YOU DO NOT USE THE DEFAULT PAYROLL SYSTEM
.sp
.li
PAY INTERVAL USAGE
.in5
.ti-5
SEMI-MONTHLY [Y] END OF LAST SEMI-MONTHLY PAY PERIOD [12/31/79]:
If any employees are paid on a semi-monthly schedule, you would state
that the semi-monthly interval is to be used, and enter the last day
of the previous semi-monthly pay period, either the 15th of the month or
the last day of the month.
.sp
.ti-5
MONTHLY [N] END OF LAST MONTHLY PAY PERIOD [12/31/79]:	If any employees
are paid once a month, you would state that the monthly interval is to
be used, and enter the last day of the previous month.
.sp
.ti-5
BIWEEKLY [N] END OF LAST BIWEEKLY PAY PERIOD [12/31/79]: If any employees
are paid every two weeks, you would state that the biweekly interval is
to be used, and enter the last day of the previous biweekly pay period.
This can be any day you want, although if both the bi-weekly and the weekly
intervals are used, the last day of the last period should be entered so
they are the same every two weeks.
.sp
.ti-5
WEEKLY [N] END OF LAST WEEKLY PAY PERIOD [12/31/79]: If any employees are
paid weekly, you would state that the weekly interval is to be used, and
enter the last day of the previous weekly pay period.
.qi
.sp2
DEFAULT EMPLOYEE VALUES: Default values are values used in the absence of other
instructions or information.  The Default Employee Values are the values
that will appear as the current answers to the requests issued by the three
employee entry programs (see Chapter 11).
They are intended to speed system setup
by making it unnecessary to type in the same answers over and over for
many employees.  If, for example, most employees accumulate vacation at the
same rate each pay period, you would enter that rate as the default.  It
is faster to change the vacation rate for the few exceptions than to enter
the same value for the many who accumulate vacation at the same rate.
.in5
.sp
.ti-5
VACATION RATE [xxxxx0.42]: This is the default rate in days or hours
(use one or
the other consistently to avoid confusion) an employee accumulates vacation
each pay period.  For example, if an employee paid semi-monthly
receives 10 days
per year, he or she would accumulate vacation at a rate of .42 days per
period (10 days/24 pay periods = .42 days per month).
.sp
.ti-5
SICK LEAVE RATE [xxxxx0.21]: This is the default rate in days or hours at which
an employee
accumulates sick leave each pay period.  For example, if an employee paid
on a semi-monthly schedule receives 5 days per year, he or she would
accumulate sick leave at a rate of .21 days per pay period (5 days/24
pay periods = .21 days per period).
.sp
.ti-5
RATE1 DESCRIPTIVE NAME [REGULARxxxxxxxx] AND DEFAULT VALUE [xxxxx1.00]:
This is the first of four user-definable pay rate types.  The four rates
usually take the descriptive names shown below (these are the default
descriptive names), but you can change them to meet your particular needs.
.sp
.li
	     RATE1 -- Regular Pay
	     RATE2 -- Overtime Pay
	     RATE3 -- Special Overtime
	     RATE4 -- Commission Pay
.sp
Although the RATE1 through RATE4 pay rates can be set for each employee
individually, the RATE2 and RATE3 default values are multiples of the
default RATE1 amount.  For example, default RATE2 can
be established as 1.5 times
RATE1, which would correspond to the usual practice of paying employees
one and one-half times their regular rate for overtime or other special
working conditions.
.sp
.ti-5
RATE2 DESCRIPTIVE NAME [OVERTIMExxxxxxx] AND FACTOR [xxxxx1.50]: The
RATE2 default value is the result of multiplying the RATE1 default amount
times the RATE2 FACTOR.  For example, if the RATE1 default value is $5.00
per hour, and the RATE2 FACTOR is 1.5, the RATE2 default value would be
$7.50 per hour.  Of course, you can change the RATE2 value for any employee.
.sp
.ti-5
RATE3 DESCRIPTIVE NAME [SPEC. OVERTIMEx] AND FACTOR [xxxxx2.00]: The RATE3
default value is the result of multiplying the RATE1 default amount times
the RATE3 FACTOR.  For example, if the RATE1 default value is $5.00 per hour,
and the RATE3 FACTOR is 2.0, the RATE3 default value would be $10.00
per hour.  Of course, you can change the RATE3 value for any employee.
.sp
.ti-5
RATE4 DESCRIPTIVE NAME [COMMISSIONxxxxx] AND EXCLUSION TYPE [1]: RATE1,
RATE2, and RATE3 pay transactions are subject to all taxes and deductions
that apply to the employee for whom they are entered.  The tax status
of RATE4 transactions, however, can be defined by the user for each employee
according to a
tax exclusion code (see Chapter 20). The Exclusion Type you enter here is
a default value that can be changed for any employee.
Exclusion type 1 means Rate4 earnings
are subject to all taxes and deductions.
.sp
.ti-5
DEFAULT PAY INTERVAL [S]: Enter the pay schedule by which most employees
are paid.  This default value is intended to speed the Employee Rate
Entry setup step.  If most employees are paid on a semi-monthly schedule
you would give that as the Default Pay Interval.  The pay interval for
employees not paid semi-monthly would have to be changed to the appropriate
interval individually.
.sp
.ti-5
DEFAULT PAY TYPE (HOURLY OR SALARIED) [S]: If most employees are salaried,
you would set the default pay type to Salaried.  Again, this is intended
as a convenience.  The Pay Type you enter for an employee only affects
certain report printing procedures.  Don't worry if an employee does not
fit neatly into one or the other classification.
.sp
.ti-5
MAX PAY FACTOR [xxxxx3.00]: The number you enter at this request, multiplied
by the RATE1 amount produces the maximum amount of pay an employee can
receive without an error message being issued when the Apply Payroll and
Print Payroll Journal programs are run.
It serves as a check to make sure no employee is paid an
inordinate amount (due to any kind of error) in a given pay period.  For
example, entering 2.0 as the MAX PAY FACTOR would ensure that any employee
paid over twice the normal amount would be brought to the attention of the
operator.
The maximum amount a particular employee can receive appears in the MAX PAY
field in the Employee Rate Entry program, and of course can be changed
for any employee.
.sp
.ti-5
NORMAL UNITS [xxxxx0.00]: This is the default value for the normal number of
pay units an employee accumulates in one pay period.  Most hourly employees
who are paid weekly work 40 hours a week, and thus would have 40 normal
units.	The same workers, if paid every two weeks would have 80 normal
units.	The normal number of pay units for salaried employees is usually 1
(one) per period.  This parameter value is only a default and is intended
to speed up the entry of employee pay rate information in the Employee Pay
Rate Entry program.  It can be changed for any employee.
.pp
A special RATE0 transaction can be entered through the Pay Transaction
Entry program that causes the Payroll System to calculate any pay units
over the Normal Units figure at RATE2.	Units up to the Normal Units amount
are paid at RATE1.  Thus, if an hourly employee normally works a 40 hour
week, and is paid at RATE2 for hours worked over the normal 40 hours,
a RATE0 transaction for 48 units would cause the employee to be
paid at RATE1 for the first 40 hours, and RATE2 (the overtime rate) for
the last 8 hours.
.qi
.sp2
DISK DRIVE ALLOCATION: The data and parameter files created and used by
the Payroll System vary in size depending on several factors.  Some files
store information used to control payroll processing (such as the Parameter,
Tax Table, Account, and System Deduction files), or summarized
historical information for printing government forms or Payroll System
reports (such as the Company History or Employee History files).  These
files (except for the Employee History file)
are generally quite small.  Other files, such as the Employee Master,
Pay Transaction Batch, Pay Detail, and Pay Check files may vary in size
depending on the number of employees on the system or the number of
pay transactions per employee per pay period.  For example, if a pay
transaction is entered daily for hourly employees who are paid once
a month, the transaction batch file may grow quite large.  Each
record on the Employee Master file is over 700 bytes long and with many
employees, the file can grow large also.
.pp
During or after certain
processing steps more than one version
of a file may exist at the same time.
During the Apply Payroll step, for example, two versions of the Pay
Detail file, an old version and a new version, exist simultaneously.
When you sort the Pay Transaction batch file, both the sorted and unsorted
versions exist on the disk together until after the Apply Payroll
program has been run to completion.  An old version and a current version
of the Check file exit simultaneously also.
If your disk storage capacity is
limited and you expect disk space to be critical, you should monitor
the size of your files (using the CP/M STAT command) throughout the pay
period so you can make adjustments well before payday arrives.	If you
have a dual-drive single density microcomputer system and find yourself
cramped for disk space, you might seriously consider going to a
double density system, or adding additional disk drives. Otherwise
you may find yourself needing to pull files off the data disk to create
more space.
.pp
The Payroll System creates and expects to find two parameter files
(PYPR1.101 and PYPR2.101)
on the drive you name as the SYSTEM FILE DRIVE the first time you run the
Payroll System.  (If you choose the default Payroll System option, drive B
is assumed and no request issued.)
The SYSTEM FILE DRIVE must be either drive A or drive B.
Most users should make drive B the SYSTEM FILE DRIVE.  The CRT.DEF file,
if used, should always reside on the currently logged drive.  Most of the files
listed below can be assigned to any drive, although the vast majority
of users should assign them to drive B.  Users with more than two
disk drives may want to assign the Pay Transaction batch file or the
Employee Master and History files
to drive C or D if the number of transactions or
employees is large.
.in5
.sp
.ti-5
TAX TABLES [B]: These files contain the state or local tax tables.  Each
tax file consists of one record and is thus very small.
.sp
.ti-5
SYSTEM DED/EARN [B] (PYDED.101): This is a very
small file containing system-wide
earnings and deduction information and the Federal tax table.  It is
created by the Deduction/Earning Entry program.
.sp
.ti-5
COMPANY HISTORY [B] (PYCOH.101): This is a small file
created and updated by the
Apply Payroll program.	It holds quarter and year to date gross
pay, direct payroll expenses, and tax withholding statistics.
.sp
.ti-5
PAYROLL ACCOUNTS [B] (PYACT.101): This file holds the account
names and numbers that
correspond to the Payroll System Account Reference Numbers.  If you select
the General Ledger interface option, it cross-references the Payroll
System Account Reference Numbers with the SSG General Ledger account
numbers.  It is a small file, with one record for each account.
.sp
.ti-5
EMPLOYEE MASTER [B] (PYEMP.101): This file holds information for calculating
employee earnings and deductions, as well as descriptive information
used by the various Payroll System printing programs.  There is one
very larger record (over 700 bytes) for each employee.	The file may be
large if many employees are on the system.
.sp
.ti-5
EMPLOYEE HISTORY [B] (PYHIS.101): This file contains summarized
quarter and year to
date information on employee earnings and deductions.  It is created and
added to by the Employee Entry and Employee Deduction/Earnings and Cost
Distribution
programs.  The Apply Payroll program updates the information on the
file.  There is one record for each employee.  The file may be large
with many employees on the system. The Employee History and Employee Master
file always exist at the same time and contain the same number of records
under normal circumstances.
.sp
.ti-5
TRANSACTION BATCH [B] (PYHRS.001): This temporary file holds the batch of pay
transactions to be processed by the next run of the Apply Payroll
program.  After sorting, a sorted version (PYHRS.101)
and an unsorted version (PYHRS.001) of the
file exist on the same disk.
The Payroll Application program completely consumes the sorted version, and
erases the unsorted version at the end of processing.
The file may be large or small, depending on the number of
employees and the number of transactions per pay period.
.sp
.ti-5
PAY CHECK [B] (PYCHK.101): The Pay Check file is
created by the Apply Payroll
program.  Each record contains only the information needed to produce a
paycheck for each employee to be paid for the period.  The Apply Payroll
program creates a backup version of the Pay Check
file (PYCHK.102) before creating a new one.
.sp
.ti-5
PAY DETAIL [B] (PYPAY.101): The Pay Detail file
contains detailed pay transaction
information for employees who are to be paid on the next run of the
Apply Payroll program for their pay interval,
and employees who were paid on the last
payroll application.  Once created the Pay Detail file is permanent.
The Apply Payroll program creates a
backup version of the old Pay Detail file before it creates the new one,
and erases it when the program ends.
The file may be
large or small depending on the number of employees and the number of
transactions per employee per pay period.
.qi
.sp2
.li
GENERAL LEDGER
.in5
.ti-5
GL INTERFACE USED (Y OR N) [N]:  This request tells the Payroll System
if it should create a file of General Ledger transactions,
and require the GL data disk with the chart of
accounts file to be "on-line" (in one of the disk drives)
when you run the Account Entry program.
.sp
.ti-5
GL BATCH FILE OUTPUT DRIVE [B]: This is the drive location where the
Payroll Application program should create the General Ledger batch file.
.sp
.ti-5
GL COA FILE SUFFIX [SS]: This is the two character suffix that identifies
the GL chart of accounts file as that of a particular company.	See you
GL manual for details.
.sp
.ti-5
GL DATA DRIVE [A]: This is the drive into which the GL data disk with the chart
of accounts should be inserted when you run the Account Entry
program, if you have selected the GL interface option.	It should never be
the same drive as the Account file or the Parameter files (the System File
Drive).
.qi
.sp2
.li
USER DEFINED PROGRAM
.in5
.ti-5
PROGRAM USED [N]: The Payroll System menu provides room for one program
created by the user.  To take advantage of this feature you would answer
Yes to this request.
.sp
.ti-5
PROGRAM [xxxxxxxx]: Give the eight character filename of the program you
created.
.sp
.ti-5
DESC. [xxxxxxxxxxxxxxxxxxxxxxxxxxxxxx]: Space is provided for a brief
description of the program. The description you enter here appears
on the Payroll System menu.
.qi
.sp2
.li
CHECK WRITING
.in5
.ti-5
COMPUTER CHECK PRINT (Y OR N) [Y]: To tell the Payroll system you want
to use the automatic check printing feature, you would answer Yes to this
request.
.sp
.ti-5
VOID CHECKS OVER MAX (Y OR N) [N]: As an added safety feature, you can
order any checks written over a specified maximum amount to be voided.
.sp
.ti-5
MAXIMUM AMOUNT [xxxx5,000.00]:	This is the most a check can be written
for without being voided if you answer the previous request affirmatively.
.qi
.sp2
.li
PAYROLL SYSTEM ACCOUNTS
.in5
.ti-5
OFFSET CASH ACCOUNT [xx1]: This is the Payroll System account reference
number to be credited with the amount of the payroll checks written.
You cannot change this account reference number
unless the account to which it is
to be changed already exists on the account file.  If you do not choose
to use the automatic check writing feature, you might want to change this
to an accounts payable account.
.sp
.ti-5
DEFAULT DIST ACCOUNT [xx3]: This is the Payroll System account reference
number debited for 100% of an employee's total cost
to the company (gross pay plus company payroll expenses)
in the absence of other
accounts distribution information for the employee.  You cannot change this
account number unless the account to which it is to be changed already
exists on the account file.
.sp
.ti-5
DISTRIBUTE EMPLOYEE COST TO 1-5 GL ACCOUNTS (Y OR N) [N]: If you want to
be able to distribute the total cost of an employee to more than one
account, answer Yes.
.qi
.sp2
.ce
SORT PARAMETERS
.pp
The Pay Transaction file is the only file in the Payroll System that is
sorted before processing (specifically, for the Apply Payroll
program).  A set of sort parameters (input drive, workfile drive, sorted
output drive location)
must be created a some point for use by the sorting program.
The input drive is assumed to be the drive location specified for the
Pay Transaction file during Parameter Entry.  The sorting program creates the
temporary
work files, and the sorted output file
on the same drive as the input file, and consequently
there are no requests
for this information.  Whenever the Pay Transaction file drive location
is changed in Parameter Entry, you must create a new sort parameter file.
At first time startup, the system creates one automatically after the
Account Entry and Deduction/Earning Entry programs have been run.  After
that, you choose the Create Sort Parameters program from the system menu
to create a new sort parameter file.  The sort parameter file (PYHOURS.SRT)
is created on the System File drive.
.he
